\section{Electroencefalografía y análisis de señales}

\subsubsection{Fundamentos del EEG}
\subsubsection{Canales, referencias y montajes}
\subsubsection{Artefactos en EEG}
\subsubsection{Crisis epilepticas y su manifestación en EEG}
\subsubsection{Segmentación temporal y anotaciones}

\section{Representación y extracción de características del EEG}
\subsection{Características temporales}
\subsection{Características espectrales}
\subsection{Sistemas dinámicos y linealización}
\subsubsection{Análisis de componentes independientes (ICA) (ICA)}\label{Marcoteorico_ICA}

\section{Aprendizaje automático supervisado}
\subsection{Aprendizaje automático tradicional}
\subsubsection{Máquina de vectores de soporte}
\subsubsection{XGBoost}
\subsubsection{Random Forest}
\subsection{Aprendizaje profundo}
\subsubsection{Redes convolucionales}
\subsubsection{Mecanismo de atención y Transformers}

\subsection{Interpretabilidad}
\subsubsection{Gradientes integrados}
Incrementa el paso de la entrada (gradual de 0 a 100 el valor de entrada) o sea un barrido, y en cada paso se analiza la pendiente
\subsubsection{Oclusión} \label{ssub:xai_oclussion}
Coloca un segmento en cero y vuelve a colocar el modelo para analizar la importancia del segmento en la respuesta.
\subsubsection{LIME}